\chapter{Discussion and Future Work}
\label{ch:conclusions}

\section{What was established}

This thesis set out to measure whether a quantized spiking neural network
reasons like its full-precision parent, and found that the measurement could not
be interpreted until the instrument was calibrated. Building and calibrating that
instrument is what the work delivers.

A four-block convolutional spiking network classifies four families of transient
noise in LIGO data at $0.9833 \pm 0.0014$ validation accuracy over three seeds,
under constraints --- shift-implementable membrane decay, at most eight time
steps, 112-pixel input --- that keep the configuration deployable by
construction rather than by subsequent rounding. Each of those constraints was
shown to cost nothing measurable in accuracy, and the resolution choice was shown
to improve the health of the deepest layer by 12.7 percentage points of dead
fraction.

The Spike Activation Map was implemented, its summation range corrected, and its
free parameter fixed on the temporal axis it governs rather than on
time-aggregated maps that cannot see it. Its faithfulness was tested
quantitatively and found to hold on two classes of four.

The central result is the calibration. Three levels were established: the
divergence produced by weight noise as large as an entire 8-bit quantization step
(PCC $0.975$), by a thirtyfold change in the explainer (PCC $0.960$), and by
initialization alone (PCC $0.153$). The third is far larger than expected, and it
carries three consequences.

\paragraph{Spatial metrics cannot compare independently trained models.} Two
networks that agree to within one validation sample in accuracy produce
explanations correlating at $0.153$. The mechanism is representational
permutation rather than any defect of the metric, and the disagreement is uniform
across samples rather than driven by outliers.

\paragraph{They can compare a model with its own parent.} Under a perturbation 85
times smaller than the 2-bit quantization step, PCC remains above $0.975$. The
two results answer different questions, and the distinction must be made
explicitly whenever such a number is reported.

\paragraph{Published figures are not interpretable without these floors.} Values
around $0.85$ are presented in the literature as evidence of preserved
explanations. For this architecture that value falls between the numerical floor
and the independent-solutions level, so it cannot be read without knowing which
comparison produced it. Establishing the floors requires no training runs beyond
a multi-seed baseline that is good practice regardless.

A fourth observation cuts across all of them: the class sensitivity ordering is
identical under three unrelated interventions. Explanation sensitivity is a
property of the class morphology, not of the perturbation applied, and
Extremely\_Loud --- saturated and spatially extensive --- registers no change
under any measure tested.

\section{Limitations}

\paragraph{Generalization of the thresholds.} The three calibration levels are
specific to this architecture, this dataset and this explainer. The
\emph{mechanism} behind the collapse, representational permutation between
independently trained networks, is general and well documented; the
\emph{magnitude} is not, and a different architecture could place the floors
elsewhere. What generalizes is the claim that the floors must be measured, not
the values measured here.

\paragraph{Labels are model output.} Accuracy is agreement with the Gravity Spy
convolutional classifier rather than with physical ground truth. A model that
reproduced that classifier's systematic errors would score well here. The
per-class confidences are available and would permit stratifying accuracy by the
reference classifier's own certainty, which was not done.

\paragraph{The partition is random with respect to time.} If glitches of the same
class are correlated within observing periods, a random test set overestimates
generalization to new periods. A block split by observing run would be the more
conservative validation; the timestamps were dropped during feature selection for
the reasons given in Section~3.1.3 and would have to be recovered.

\paragraph{$T$ sits at the imposed ceiling.} The search was forbidden from
exploring more than eight time steps, so it is not known whether it would have
preferred a longer integration window. An unconstrained search selected seven,
which bounds the likely cost but does not eliminate the question.

\paragraph{The deletion metric fails on two classes.} The geometric explanation
offered in Section~5.4.2 --- that deletion-based tests penalize spatially
extended evidence --- is an interpretation, not a demonstration. An insertion
curve, which rewards rather than penalizes distributed evidence, would test it
directly.

\paragraph{The explanation is class-agnostic.} SAM does not use labels, so the
deletion metric measures faithfulness to the model's \emph{predicted} class. This
is the correct reading of the metric and matches its original formulation, but it
means misclassified examples are informative and should be inspected separately
rather than pooled.

\section{The quantization study}
\label{sec:grid}

The calibrated instrument exists to support a specific experiment, specified here
in full.

\subsection{The grid}

Quantization-aware training is implemented by fake-quantization with a
straight-through estimator~\cite{bengio2013ste}. The critical design choice is to
\textbf{quantize the two targets separately before quantizing them together},
which converts a single observation into a causal ablation and is what makes the
hypothesis of Chapter~\ref{ch:introduction} testable.

\begin{table}[htbp]
\centering
\begin{tabular}{llll}
\toprule
Arm & Weights & Membrane & Purpose \\
\midrule
FP32 reference & 32-bit float & 32-bit float & baseline for every comparison \\
Weight-only & 8 / 4 / 2 bit & 32-bit float & isolates the synaptic effect \\
Membrane-only & 32-bit float & 8 / 4 / 2 bit & isolates the state effect \\
Joint (MINT-style) & 8 / 4 / 2 bit & 8 / 4 / 2 bit & deployable, shared scale \\
\bottomrule
\end{tabular}
\caption{Ten configurations, each at three seeds.}
\label{tab:grid}
\end{table}

The two arms have \emph{opposite} sensitivities to the resolution choice, which
independently justifies treating them as separate ablations. Because $270\,240$
of the parameters are convolutional and therefore resolution-independent, the
weight arm achieves nearly the same absolute memory reduction at any resolution,
whereas the membrane state falls twelvefold between 224 and 64 pixels.

\subsection{Divergence metrics and what they can claim}

For every configuration, SAM maps are recomputed and compared against the
reference at the same seed, using: spatial Pearson correlation, top-20\,\%
intersection-over-union, structural similarity~\cite{wang2004ssim}, the
per-timestep correlation curve, the temporal centre of mass, cross-faithfulness,
and the Earth Mover's Distance between firing-rate distributions for
comparability with~\cite{smith2026}.

Three constraints follow directly from Chapter~\ref{ch:results} and must be
applied when those numbers are read.

Results below the numerical floor of Table~\ref{tab:calibration} are measurement
noise. IoU saturates early and is to be treated as a sensitive early indicator
rather than a graded measure. And \textbf{cross-faithfulness will be
interpretable on Blip and Extremely\_Loud only}, since the deletion metric does
not separate SAM from a random baseline on the other two classes, leaving no
baseline signal for a quantized model to degrade.

\subsection{Stratified reporting}

An aggregate divergence per configuration admits only a yes-or-no answer. Joining
the per-sample outputs to the physical descriptors allows the central result to be
reported as a relation instead.

The hypothesis is that divergence increases as event strength decreases, with a
mechanism specific to spiking networks: quantizing the membrane potential flips a
neuron's spike decision only where the potential lies within the quantization
error of the threshold, and weak events produce potentials that sit close to
threshold rather than overshooting it. Because a flipped spike is a discrete
event that propagates both to the next layer and forward in time through the
decay and reset, and because SAM is constructed from spike trains, the flip
changes the object the explanation is made of rather than merely perturbing it.

The prediction discriminates between arms: membrane quantization applies a fixed
absolute error and should show a steep dependence on event strength, whereas
weight quantization perturbs the input current approximately in proportion to the
signal and should be flatter.

Two preconditions were verified. SNR does translate into image intensity, the
Gravity Spy colour scale being fixed rather than auto-scaled per image. And SNR
is very nearly a proxy for class --- the top global SNR quartile is 92.8\,\%
Extremely\_Loud --- so stratification must use \textbf{within-class quantiles},
or any trend measured would be a class effect wearing a strength label.

\subsection{Energy accounting}

Synaptic operations are counted as
\begin{equation}
\mathrm{SOP}_{\ell} = f_{\ell-1} \cdot T \cdot \mathrm{MAC}_{\ell},
\end{equation}
where $f_{\ell-1}$ is the firing rate of the preceding layer and
$\mathrm{MAC}_{\ell}$ the multiply-accumulate count of the equivalent dense
layer, and multiplied by the per-operation costs of~\cite{horowitz2014}. The
rates must be the per-hidden-layer figures measured on the converged checkpoint.

\textbf{One layer is asymmetric and dominates the budget.} Under direct encoding
the first convolution receives analog pixel values rather than spikes, so its
operations are expensive multiply-accumulates. Because the input is identical at
every time step, the current it produces is constant in time and is computed once
rather than $T$ times. This is a standard implementation of constant encoding,
but stating it explicitly matters: recomputing that layer at every step would
make the spiking network \emph{more} expensive than its convolutional
equivalent, and the two readings differ by a factor of $2.9$.

The headline figure then plots accuracy, explanation fidelity and energy on a
shared bit-width axis. The three curves falling together would be unremarkable.
The result to look for is a \textbf{divergence}: an operating point where
accuracy is still intact and energy already much reduced, but the explanation has
quietly stopped describing the same reasoning. Anyone selecting a configuration
on accuracy alone would choose that point and deploy a model that decides
correctly for different reasons. Should the curves instead move together, that
too is reportable: it would establish that for this architecture accuracy is an
adequate proxy for explanation fidelity under compression, which is the
assumption the literature currently makes without testing.

\section{Closing remark}

The question this thesis asked has not been answered. What it established instead
is the condition under which the answer will mean something --- and, along the
way, that a class of published results in this area rests on a measurement whose
noise floor was never checked. For a field in which compressed models are
increasingly deployed where their explanations are acted upon, that seems worth
knowing before the next number is reported.
